\documentclass[10pt,a4paper]{amsart}
\usepackage[margin=19mm]{geometry}
\usepackage{amsmath,amssymb,mathtools}
\usepackage[scr=rsfs,cal=euler]{mathalfa}
\usepackage{xfrac}
\usepackage[unicode,hidelinks,bookmarks=false]{hyperref}
\newcommand{\ie}{i.e.\ }
\newcommand{\cf}{cf.\ }
\newcommand{\resp}{respectively\ }
\newcommand{\Subsection}{\subsection}
\providecommand{\nobreakdash}{\nobreak\hbox{-}}
\setlength{\emergencystretch}{2em}
\multlinegap0pt
\usepackage{bm}
\usepackage{bbm}
\usepackage{dsfont}
\usepackage{xspace}
\usepackage{cancel}
\usepackage{mathtools}
\usepackage{amsthm}
\makeatletter
\@ifundefined{theorem}{\newtheorem{theorem}{Theorem}}{}
\@ifundefined{proposition}{\newtheorem{proposition}[theorem]{Proposition}}{}
\makeatother
\title[Algebraic hyperbolicity of very general hypersurfaces in wei]{Algebraic hyperbolicity of very general hypersurfaces in weighted projective spaces}
\author{Jiahe Wang}
\date{Source: arXiv:2511.04890v1 (CC BY 4.0)}
\begin{document}
\maketitle

\noindent\emph{Source and licence.} Algebraic hyperbolicity of very general hypersurfaces in weighted projective spaces. Version arXiv:2511.04890v1 (CC BY 4.0), distributed under the Creative Commons Attribution 4.0 International licence, as recorded in the arXiv licence metadata for this exact version and shown on the abstract page https://arxiv.org/abs/2511.04890v1. Full rights notice: licenses/arxiv-2511.04890-CC-BY-4.0.txt.\par\smallskip

\noindent\textit{Editorial scope.} Counted unit: the Proposition whose source label is \texttt{prop4} in the article (source0.tex lines 400--407, source label \texttt{prop4}), delivered together with the article's own complete original proof of it (source0.tex lines 410--429, one \texttt{proof} environment of 20 lines). No mathematics is written, repaired or re-derived: statement and proof are the article's own text line for line. Required context retained uncounted: prop2 (lines 346--348); prop3 (lines 383--385). Every definition, display and label of those lines is retained unchanged. Omitted from this package and not redistributed: 1--345, 349--382, 386--399, 408--409, 430--476. Nothing inside a retained excerpt is omitted or altered: every retained body is delivered exactly as the article's lines are written, so no omission receipt of any kind accompanies this package. Citations: the retained excerpts carry no citation command at all, so no bibliography entry of the article is transcribed and no reference is redistributed. Reference adaptation: none was needed and none was made; every reference inside the delivered unit resolves inside it, so no unresolved reference or citation is produced. Printed slips kept literal: no printed word, symbol, hypothesis or number of the counted statement or of its proof was altered. Layout and class: the article's own class is \texttt{amsart} with packages including amsmath, amssymb, tikz-cd, array, float, amsthm, enumitem, mathrsfs, thmtools, thm-restate, hyperref, cleveref; the wrapper is the lane's pinned \texttt{amsart} editorial wrapper at 10pt on A4, which restates the theorem-like environments the retained text needs and adds the article's own macro definitions that the retained text needs (none), with extra wrapper packages (bm, bbm, dsfont, xspace, cancel, mathtools). The delivered numbering is the unit's own. Assets: no retained span contains a figure environment, a graphics-inclusion command, a PDF-page inclusion, a table or a TikZ picture; no image file is copied into this package, the package declares no asset dependency and no image file is redistributed with it. Licence: version arXiv:2511.04890v1 of this article is distributed under the Creative Commons Attribution 4.0 International licence, as recorded in the arXiv licence metadata for this exact version and shown on the abstract page https://arxiv.org/abs/2511.04890v1. A full rights notice is delivered at licenses/arxiv-2511.04890-CC-BY-4.0.txt. Characters: every retained excerpt is pure ASCII, so the delivered engine, XeLaTeX with its default Latin Modern text font, reports no missing character.
\begin{proposition}\label{prop2}
 If $m > \frac{a_0 + \dots + a_n}{a_0a_1... a_n} + (n-2)$, then a very general hypersurface $X$ of $|mH|$ is algebraically hyperbolic outside the toric boundary.
\end{proposition}

\begin{proposition}\label{prop3}
For a weighted projective $3$-fold $\mathbb{P}(a_0,a_1,a_2,a_3)$ with isolated singularities, if $\mathbb{P}(a_0,a_1,a_2,a_3) \neq \mathbb{P}(1,1,1,n)$ or $\mathbb{P}(1,1,2,3)$, then a very general surface $X\in |mH|$ is algebraically hyperbolic if $m \geq 2$. For $\mathbb{P}(1,1,2,3)$, a very general surface $X\in|mH|$ is algebraically hyperbolic if $m\geq 3$.
\end{proposition}

\begin{proposition}\label{prop4}
  
Let $\mathbb{P} = \mathbb{P}(a_0,\dots,a_n)$ be a weighted projective space with isolated singularities. Let
\[
\Theta := \max_{\substack{I \subset \{0,\dots,n\} \\ |I|\ge 4}} \left\{\frac1{\prod_{i\notin I}a_i} \left( \frac{\sum_{i \in I} a_i}{\prod_{i\in I} a_i} + (|I|-3) \right)\right\}.
\]
Then for every $m > \Theta$, a very general hypersurface $X \in |mH|$ is algebraically hyperbolic.
\end{proposition}

\begin{proof}
Since we work with $\mathbb P$ and its toric strata, we consider a very general $X$ of $mH$ in $\widetilde{\mathbb P}$ to be a very general $X$ of $\mathcal O(mk)$ in $\mathbb P$. We note that for a weighted projective space $\mathbb{P}(a_0,\dots,a_n)$, the toric boundary divisors are exactly 
\[
D_i = \mathbb{P}(a_0,\dots,\widehat{a_i},\dots,a_n), \quad i \in \{0,\dots,n\}.
\] 
Let $I \subset \{0,\dots,n\}$, $|I|\geq 4$, then we have $\mathbb{P}(a_i, i \in I)$ as a stratum of $\mathbb{P}(a_0,\dots,a_n)$ in the sense of hyperbolicity. That is, if we show that for each $I \subset \{0,\dots,n\}$ with $|I|\geq 5$, a very general section of $\mathcal O(mk)|_{\mathbb{P}_I}$ in $\mathbb{P}_I$ is algebraically hyperbolic with respect to the ample $\mathcal O(k)|_{\mathbb{P}_I}$ outside the toric boundary of $\mathbb{P}_I$, and for $|I|=4$ a very general section of $\mathcal O(mk)|_{\mathbb{P}_I}$ in $\mathbb{P}_I$ is algebraically hyperbolic with respect to $\mathcal O(k)|_{\mathbb P_I}$. Then since in our case, the restriction maps of the global sections of $\mathcal O(mk)$ to each stratum $H^{0}\!\big(\mathbb{P},\, \mathcal{O}_{\mathbb{P}}(mk)\big)
\rightarrow
H^{0}\!\big(\mathbb{P}_{I},\, \mathcal{O}_{\mathbb{P}_{I}}(mk)\big)$ are surjective. So taking the finite intersection of preimages of very general sets one in each $H^{0}\!\big(\mathbb{P}_{I},\, \mathcal{O}_{\mathbb{P}_{I}}(mk)\big)$, we get a very general set in $H^{0}\!\big(\mathbb{P},\, \mathcal{O}_{\mathbb{P}}(mk)\big)$, such that for any hypersurface $X$ in the set, $\mathbb P_I \cap X$ is algebraically hyperbolic for each $|I| = 4$ and $\mathbb P_I \cap X$ is algebraically hyperbolic outside the toric boundary for each $|I| \geq 5$, so together they imply that $X$ is algebraically hyperbolic. We note that each stratum $\mathbb P_I$ is a weighted projective space with isolated singularities, so we may apply Proposition \ref{prop2}, \ref{prop3}. To find an $\epsilon$, we take the minimum $\epsilon_I$ of each stratum. We note that algebraic hyperbolicity is independent of ample divisor, but here we fix the ample divisor $\mathcal O(k)$ to give a global $\epsilon$ for a very general $X$. 

We see that the inequality $m>\Theta$ produces a sufficient condition for algebraic hyperbolicity. This is because by Proposition \ref{prop3},
\[
m > \frac{1}{\prod_{i\notin I}a_i}(\frac{\sum_{i \in I} a_i}{\prod_{i\in I} a_i} + (|I|-3)), \quad \text{for } |I|=4,
\] 
means that a very general section of $mq_I\mathcal O(k_I)  = \mathcal O(mk)|_{\mathbb{P}_I}$, where $\mathcal O(k_I):= \mathcal O(\prod_{i\in I} a_i)$ and $q_I := \prod_{i \notin I} a_i$, is algebraically hyperbolic with respect to $\mathcal O(k_I)$. So this implies a very general section of $\mathcal O(mk)|_{\mathbb{P}_I}$ is algebraically hyperbolic with respect to $\mathcal O(k)|_{\mathbb{P}_I}$ with $\epsilon_I = \frac 1{q_I}(\frac{\sum_{i \in I} a_i}{\prod_{i\in I} a_i} + (|I|-3))$. Similarly, for $|I|\geq5$ by Proposition \ref{prop2},
\[
m >  \frac{1}{\prod_{i\notin I}a_i}(\frac{\sum_{i \in I} a_i}{\prod_{i\in I} a_i} + (|I|-3)),
\] 
implies that a very general section of $mq_I\mathcal O(k_I)  = \mathcal O(mk)|_{\mathbb{P}_I}$ is algebraically outside the toric boundary of $\mathbb P_I$ with respect to $\mathcal O(k)|_{\mathbb{P}_I}$ with $\epsilon_I = \frac 1{q_I}(\frac{\sum_{i \in I} a_i}{\prod_{i\in I} a_i} + (|I|-3))$. So we get a very general section of $\mathcal O(mk)$ is algebraically hyperbolic with respect to $\mathcal O(k)$ with an $\epsilon = \operatorname{min}\{\epsilon_I, |I|\geq4\}$.
   
\end{proof}

\end{document}
